% ============================================================
%  Related Work — % ============================================================
%  Related Work — % ============================================================
%  Related Work — % ============================================================
%  Related Work — \input{2_related_work} in main paper
% ============================================================

\section{Related Work}
\label{sec:related-work}

\paragraph{RAG over long context.}
Retrieval-Augmented Generation~\cite{rag_2020} is an efficient way to extend
LLM reasoning to out-of-distribution data, but its performance degrades in
long-context settings: recall slows as context grows and accuracy can
actively decrease, suggesting a return of hallucination~\cite{leng_long_2024}. Graph-based RAG methods address this through richer offline indexing.
HippoRAG~2~\cite{hipporagv2} combines knowledge-graph triples with raw passages in a unified index and uses an LLM-driven recognition step plus Personalized PageRank to surface
relevant evidence; it improves on earlier structured architectures such as
LightRAG~\cite{lightrag} and GraphRAG~\cite{graphrag_2024}, but its heavy
LLM-driven indexing complicates incremental deployment, and chunking remains
a fixed-size, query-agnostic preprocessing step. These methods are also token-intensive, especially during graph construction. LinearRAG~\cite{linearrag_2025} has been proposed to reduce time and token consumption by scaling linearly with graph size, and shows that graph-based RAG often introduces relational noise that hurts generation. Closer to our work, Adaptive-RAG~\cite{adaptive_rag_2024} routes
queries to retrieval strategies of varying complexity based on predicted
question difficulty.
Our query-routing contribution extends this intuition to the
\emph{multi-source} setting, conditioning retrieval on predicted source
relevance rather than complexity.

\paragraph{Multi-hop question answering.}
Multi-hop QA requires aggregating evidence from multiple documents.
Benchmarks such as HotpotQA~\cite{hotpotqa_2018} and
MuSiQue~\cite{musique_2022} have driven significant progress, but they rely
on clean modern text and assume evidence resides within a single,
homogeneous corpus.
Historical QA introduces compounding difficulties: OCR noise, temporal and
orthographic drift, and the need to bridge
\emph{heterogeneous} source types whose retrieval dynamics differ
substantially.
Standard multi-hop benchmarks provide no guidance for this cross-corpus
setting.


% Là tu parles de sources parlementaires, mais on en parle pas du tout avant donc je remplacerais par 

\paragraph{RAG for historical archives.}
Recent work highlights how OCR noise, orthographic
variation, and semantic drift limit standard retrieval pipelines.
\citet{Mudet_Bakkali_2025} stabilise recall on noisy historical corpora via
semantic query expansion and Reciprocal Rank Fusion, but their evaluation
is restricted to documents of at most 512 tokens and to a single source
type --- leaving the cross-source, long-context regime of historical
archives largely unexplored.

% Il faudrait ajouter des référebces pour "Recent work", non ?
% Pour moi : soit tu ajoutes des références, soit tu fusionnes avec la première section ("RAG et contexte long")

%\paragraph{RAG and history.}
%Recent work highlights how OCR noise, orthographic variation, and semantic drift limit standard retrieval pipelines. \citet{Mudet_Bakkali_2025} stabilise recall on noisy historical corpora via semantic query expansion and Reciprocal Rank Fusion, but their evaluation is restricted to documents of at most 512 tokens and to a single source type -- leaving the cross-source, long-context regime of parliamentary archives largely unexplored.


\paragraph{Segmentation in RAG.}
Segmentation is the first step of RAG indexing and has a direct effect on
retrieval quality.
Classical fixed-size segmentation makes a one-size-fits-all assumption;
recursive and LLM-guided alternatives have been
proposed~\cite{chromadb_chunking_2024}.
\citet{hope_2025} introduce a domain-agnostic framework for evaluating
chunking strategies and show that semantic independence between passages is
the most consequential factor for downstream performance.
Recent works have highlighted the importance of moving away from the one-size-fits-all hypothesis and have studied the importance of adapting the chunking method to the domain of study~\cite{reconstructingcontextevaluatingadvanced,mocmixturestextchunking}.
Following those ideas, approaches have used document-specific chunking optimisation research showing how document-specific chunking optimisation can improve RAG results~\cite{adaptive_chunking}. These approaches still treat the chunking phase and the query phase as independent, focusing on off-line optimisation. We argue that domain experts pose heterogeneous questions, and that the optimal segmentation depends both on the document and on the query.
Segmentation should therefore be treated as a query-conditioned parameter, jointly optimised with retrieval, rather than as a fixed pre-processing choice. in main paper
% ============================================================

\section{Related Work}
\label{sec:related-work}

\paragraph{RAG over long context.}
Retrieval-Augmented Generation~\cite{rag_2020} is an efficient way to extend
LLM reasoning to out-of-distribution data, but its performance degrades in
long-context settings: recall slows as context grows and accuracy can
actively decrease, suggesting a return of hallucination~\cite{leng_long_2024}. Graph-based RAG methods address this through richer offline indexing.
HippoRAG~2~\cite{hipporagv2} combines knowledge-graph triples with raw passages in a unified index and uses an LLM-driven recognition step plus Personalized PageRank to surface
relevant evidence; it improves on earlier structured architectures such as
LightRAG~\cite{lightrag} and GraphRAG~\cite{graphrag_2024}, but its heavy
LLM-driven indexing complicates incremental deployment, and chunking remains
a fixed-size, query-agnostic preprocessing step. These methods are also token-intensive, especially during graph construction. LinearRAG~\cite{linearrag_2025} has been proposed to reduce time and token consumption by scaling linearly with graph size, and shows that graph-based RAG often introduces relational noise that hurts generation. Closer to our work, Adaptive-RAG~\cite{adaptive_rag_2024} routes
queries to retrieval strategies of varying complexity based on predicted
question difficulty.
Our query-routing contribution extends this intuition to the
\emph{multi-source} setting, conditioning retrieval on predicted source
relevance rather than complexity.

\paragraph{Multi-hop question answering.}
Multi-hop QA requires aggregating evidence from multiple documents.
Benchmarks such as HotpotQA~\cite{hotpotqa_2018} and
MuSiQue~\cite{musique_2022} have driven significant progress, but they rely
on clean modern text and assume evidence resides within a single,
homogeneous corpus.
Historical QA introduces compounding difficulties: OCR noise, temporal and
orthographic drift, and the need to bridge
\emph{heterogeneous} source types whose retrieval dynamics differ
substantially.
Standard multi-hop benchmarks provide no guidance for this cross-corpus
setting.


% Là tu parles de sources parlementaires, mais on en parle pas du tout avant donc je remplacerais par 

\paragraph{RAG for historical archives.}
Recent work highlights how OCR noise, orthographic
variation, and semantic drift limit standard retrieval pipelines.
\citet{Mudet_Bakkali_2025} stabilise recall on noisy historical corpora via
semantic query expansion and Reciprocal Rank Fusion, but their evaluation
is restricted to documents of at most 512 tokens and to a single source
type --- leaving the cross-source, long-context regime of historical
archives largely unexplored.

% Il faudrait ajouter des référebces pour "Recent work", non ?
% Pour moi : soit tu ajoutes des références, soit tu fusionnes avec la première section ("RAG et contexte long")

%\paragraph{RAG and history.}
%Recent work highlights how OCR noise, orthographic variation, and semantic drift limit standard retrieval pipelines. \citet{Mudet_Bakkali_2025} stabilise recall on noisy historical corpora via semantic query expansion and Reciprocal Rank Fusion, but their evaluation is restricted to documents of at most 512 tokens and to a single source type -- leaving the cross-source, long-context regime of parliamentary archives largely unexplored.


\paragraph{Segmentation in RAG.}
Segmentation is the first step of RAG indexing and has a direct effect on
retrieval quality.
Classical fixed-size segmentation makes a one-size-fits-all assumption;
recursive and LLM-guided alternatives have been
proposed~\cite{chromadb_chunking_2024}.
\citet{hope_2025} introduce a domain-agnostic framework for evaluating
chunking strategies and show that semantic independence between passages is
the most consequential factor for downstream performance.
Recent works have highlighted the importance of moving away from the one-size-fits-all hypothesis and have studied the importance of adapting the chunking method to the domain of study~\cite{reconstructingcontextevaluatingadvanced,mocmixturestextchunking}.
Following those ideas, approaches have used document-specific chunking optimisation research showing how document-specific chunking optimisation can improve RAG results~\cite{adaptive_chunking}. These approaches still treat the chunking phase and the query phase as independent, focusing on off-line optimisation. We argue that domain experts pose heterogeneous questions, and that the optimal segmentation depends both on the document and on the query.
Segmentation should therefore be treated as a query-conditioned parameter, jointly optimised with retrieval, rather than as a fixed pre-processing choice. in main paper
% ============================================================

\section{Related Work}
\label{sec:related-work}

\paragraph{RAG over long context.}
Retrieval-Augmented Generation~\cite{rag_2020} is an efficient way to extend
LLM reasoning to out-of-distribution data, but its performance degrades in
long-context settings: recall slows as context grows and accuracy can
actively decrease, suggesting a return of hallucination~\cite{leng_long_2024}. Graph-based RAG methods address this through richer offline indexing.
HippoRAG~2~\cite{hipporagv2} combines knowledge-graph triples with raw passages in a unified index and uses an LLM-driven recognition step plus Personalized PageRank to surface
relevant evidence; it improves on earlier structured architectures such as
LightRAG~\cite{lightrag} and GraphRAG~\cite{graphrag_2024}, but its heavy
LLM-driven indexing complicates incremental deployment, and chunking remains
a fixed-size, query-agnostic preprocessing step. These methods are also token-intensive, especially during graph construction. LinearRAG~\cite{linearrag_2025} has been proposed to reduce time and token consumption by scaling linearly with graph size, and shows that graph-based RAG often introduces relational noise that hurts generation. Closer to our work, Adaptive-RAG~\cite{adaptive_rag_2024} routes
queries to retrieval strategies of varying complexity based on predicted
question difficulty.
Our query-routing contribution extends this intuition to the
\emph{multi-source} setting, conditioning retrieval on predicted source
relevance rather than complexity.

\paragraph{Multi-hop question answering.}
Multi-hop QA requires aggregating evidence from multiple documents.
Benchmarks such as HotpotQA~\cite{hotpotqa_2018} and
MuSiQue~\cite{musique_2022} have driven significant progress, but they rely
on clean modern text and assume evidence resides within a single,
homogeneous corpus.
Historical QA introduces compounding difficulties: OCR noise, temporal and
orthographic drift, and the need to bridge
\emph{heterogeneous} source types whose retrieval dynamics differ
substantially.
Standard multi-hop benchmarks provide no guidance for this cross-corpus
setting.


% Là tu parles de sources parlementaires, mais on en parle pas du tout avant donc je remplacerais par 

\paragraph{RAG for historical archives.}
Recent work highlights how OCR noise, orthographic
variation, and semantic drift limit standard retrieval pipelines.
\citet{Mudet_Bakkali_2025} stabilise recall on noisy historical corpora via
semantic query expansion and Reciprocal Rank Fusion, but their evaluation
is restricted to documents of at most 512 tokens and to a single source
type --- leaving the cross-source, long-context regime of historical
archives largely unexplored.

% Il faudrait ajouter des référebces pour "Recent work", non ?
% Pour moi : soit tu ajoutes des références, soit tu fusionnes avec la première section ("RAG et contexte long")

%\paragraph{RAG and history.}
%Recent work highlights how OCR noise, orthographic variation, and semantic drift limit standard retrieval pipelines. \citet{Mudet_Bakkali_2025} stabilise recall on noisy historical corpora via semantic query expansion and Reciprocal Rank Fusion, but their evaluation is restricted to documents of at most 512 tokens and to a single source type -- leaving the cross-source, long-context regime of parliamentary archives largely unexplored.


\paragraph{Segmentation in RAG.}
Segmentation is the first step of RAG indexing and has a direct effect on
retrieval quality.
Classical fixed-size segmentation makes a one-size-fits-all assumption;
recursive and LLM-guided alternatives have been
proposed~\cite{chromadb_chunking_2024}.
\citet{hope_2025} introduce a domain-agnostic framework for evaluating
chunking strategies and show that semantic independence between passages is
the most consequential factor for downstream performance.
Recent works have highlighted the importance of moving away from the one-size-fits-all hypothesis and have studied the importance of adapting the chunking method to the domain of study~\cite{reconstructingcontextevaluatingadvanced,mocmixturestextchunking}.
Following those ideas, approaches have used document-specific chunking optimisation research showing how document-specific chunking optimisation can improve RAG results~\cite{adaptive_chunking}. These approaches still treat the chunking phase and the query phase as independent, focusing on off-line optimisation. We argue that domain experts pose heterogeneous questions, and that the optimal segmentation depends both on the document and on the query.
Segmentation should therefore be treated as a query-conditioned parameter, jointly optimised with retrieval, rather than as a fixed pre-processing choice. in main paper
% ============================================================

\section{Related Work}
\label{sec:related-work}

\paragraph{RAG over long context.}
Retrieval-Augmented Generation~\cite{rag_2020} is an efficient way to extend
LLM reasoning to out-of-distribution data, but its performance degrades in
long-context settings: recall slows as context grows and accuracy can
actively decrease, suggesting a return of hallucination~\cite{leng_long_2024}. Graph-based RAG methods address this through richer offline indexing.
HippoRAG~2~\cite{hipporagv2} combines knowledge-graph triples with raw passages in a unified index and uses an LLM-driven recognition step plus Personalized PageRank to surface
relevant evidence; it improves on earlier structured architectures such as
LightRAG~\cite{lightrag} and GraphRAG~\cite{graphrag_2024}, but its heavy
LLM-driven indexing complicates incremental deployment, and chunking remains
a fixed-size, query-agnostic preprocessing step. These methods are also token-intensive, especially during graph construction. LinearRAG~\cite{linearrag_2025} has been proposed to reduce time and token consumption by scaling linearly with graph size, and shows that graph-based RAG often introduces relational noise that hurts generation. Closer to our work, Adaptive-RAG~\cite{adaptive_rag_2024} routes
queries to retrieval strategies of varying complexity based on predicted
question difficulty.
Our query-routing contribution extends this intuition to the
\emph{multi-source} setting, conditioning retrieval on predicted source
relevance rather than complexity.

\paragraph{Multi-hop question answering.}
Multi-hop QA requires aggregating evidence from multiple documents.
Benchmarks such as HotpotQA~\cite{hotpotqa_2018} and
MuSiQue~\cite{musique_2022} have driven significant progress, but they rely
on clean modern text and assume evidence resides within a single,
homogeneous corpus.
Historical QA introduces compounding difficulties: OCR noise, temporal and
orthographic drift, and the need to bridge
\emph{heterogeneous} source types whose retrieval dynamics differ
substantially.
Standard multi-hop benchmarks provide no guidance for this cross-corpus
setting.


% Là tu parles de sources parlementaires, mais on en parle pas du tout avant donc je remplacerais par 

\paragraph{RAG for historical archives.}
Recent work highlights how OCR noise, orthographic
variation, and semantic drift limit standard retrieval pipelines.
\citet{Mudet_Bakkali_2025} stabilise recall on noisy historical corpora via
semantic query expansion and Reciprocal Rank Fusion, but their evaluation
is restricted to documents of at most 512 tokens and to a single source
type --- leaving the cross-source, long-context regime of historical
archives largely unexplored.

% Il faudrait ajouter des référebces pour "Recent work", non ?
% Pour moi : soit tu ajoutes des références, soit tu fusionnes avec la première section ("RAG et contexte long")

%\paragraph{RAG and history.}
%Recent work highlights how OCR noise, orthographic variation, and semantic drift limit standard retrieval pipelines. \citet{Mudet_Bakkali_2025} stabilise recall on noisy historical corpora via semantic query expansion and Reciprocal Rank Fusion, but their evaluation is restricted to documents of at most 512 tokens and to a single source type -- leaving the cross-source, long-context regime of parliamentary archives largely unexplored.


\paragraph{Segmentation in RAG.}
Segmentation is the first step of RAG indexing and has a direct effect on
retrieval quality.
Classical fixed-size segmentation makes a one-size-fits-all assumption;
recursive and LLM-guided alternatives have been
proposed~\cite{chromadb_chunking_2024}.
\citet{hope_2025} introduce a domain-agnostic framework for evaluating
chunking strategies and show that semantic independence between passages is
the most consequential factor for downstream performance.
Recent works have highlighted the importance of moving away from the one-size-fits-all hypothesis and have studied the importance of adapting the chunking method to the domain of study~\cite{reconstructingcontextevaluatingadvanced,mocmixturestextchunking}.
Following those ideas, approaches have used document-specific chunking optimisation research showing how document-specific chunking optimisation can improve RAG results~\cite{adaptive_chunking}. These approaches still treat the chunking phase and the query phase as independent, focusing on off-line optimisation. We argue that domain experts pose heterogeneous questions, and that the optimal segmentation depends both on the document and on the query.
Segmentation should therefore be treated as a query-conditioned parameter, jointly optimised with retrieval, rather than as a fixed pre-processing choice.